We now state without proof the ``higher-order\textquotesingle\textquotesingle\ version of Lemma~\ref{lem:flat-boundary-local-regularity}.

\begin{Lemma}
\label{lem:higher-order-flat-boundary-regularity}
Let the hypotheses of Lemma~\ref{lem:flat-boundary-local-regularity} hold. In addition, assume that the coefficients satisfy $a_{ij}\in W^{k+1,\infty}(D^+)$, $b_i,c\in W^{k,\infty}(D^+)$ and that $f\in H^k(D^+)$. Then $u\in H^{k+2}(Q^+)$ and there exists a constant $C$ depending on $R$ such that
\begin{equation}\label{eq:higher-order-flat-boundary-local-regularity-estimate}
\|u\|_{H^{k+2}(Q^+)}\leq C(\|f\|_{H^k(D^+)}+\|u\|_{H^1(D^+)}).
\end{equation}

\end{Lemma}

To prove Theorem~\ref{thm:global-boundary-regularity} we must consider a general domain $\Omega$. Of course, since we did all the work of proving Lemma~\ref{lem:flat-boundary-local-regularity} on a half-space, our next task is to ``straighten out\textquotesingle\textquotesingle\ the boundary of $\Omega$. Since $\partial\Omega$ is a $C^2$ surface, we know that for each $\overline{\bm{x}}\in\partial\Omega$ there is an $R>0$ and a $C^2(\mathbb{R}^{n-1})$ function $\psi$ such that (after a possible renumbering and reorientation of coordinates)
\begin{equation}\label{eq:boundary-local-graph-coordinate-description}
\partial\Omega\cap B_R(\overline{\bm{x}})=\{\bm{x}\in B_R(\overline{\bm{x}})\mid x_n=\psi(x_1,x_2,\ldots,x_{n-1})\},
\end{equation}

\begin{equation}\label{eq:boundary-local-domain-above-graph-description}
\Omega\cap B_R(\overline{\bm{x}})=\{\bm{x}\in B_R(\overline{\bm{x}})\mid x_n>\psi(x_1,x_2,\ldots,x_{n-1})\};
\end{equation}
and moreover, the mapping
\begin{equation}\label{eq:boundary-straightening-forward-coordinate-map}
B_R(\overline{\bm{x}})\ni\bm{x}\mapsto\bm{y}=\Psi(\bm{x})\in\mathbb{R}^n
\end{equation}
defined by\[\begin{aligned}
y_i&:=x_i-\overline x_i,\quad i=1,\ldots,n-1,\\
y_n&:=x_n-\psi(x_1,\ldots,x_{n-1})
\end{aligned}\]is one-to-one. Define $\overline\Phi:=\Psi^{-1}$. Note that $\Phi$ is a $C^2$ function which transforms the set $\Omega^{\prime}:=\Omega\cap B_R(\overline{\bm{x}})$ (in what we refer to as $\bm{x}$ space) into a set $\Omega^{\prime\prime}$ in the half-space $y_n>0$ (of $\bm{y}$ space). Note also that the point $\overline{\bm{x}}$ is mapped to the origin of $\bm{y}$ space (cf. Figure~\ref{fig:boundary-straightening-original}).

\begin{figure}[H]
\centering
\includegraphics[width=.78\linewidth]{figures/chapter-09-boundary-straightening-original.png}
\caption{Straightening out the boundary.}\label{fig:boundary-straightening-original}
\end{figure}
